\problem{CM01}{正则二分图的轮次安排}{Hall 条件；完美匹配分解}
有 $n$ 名员工和 $n$ 项任务。每名员工恰能执行 $d$ 项任务，每项任务也恰有 $d$ 名员工能执行，其中 $1\le d\le n$。证明能安排 $d$ 轮，每轮每名员工执行一项不同任务，而且每条允许的员工--任务配对恰使用一次。给出所用匹配定理的证明要点。
\solution 将允许配对画成二分图。对任意员工集合 $S$，其发出的边数为 $d|S|$，这些边都到达邻集 $N(S)$；每个任务最多接收 $d$ 条，故 $d|S|\le d|N(S)|$，即 $|N(S)|\ge|S|$。
Hall 定理说，此条件保证存在覆盖全部员工的匹配。可按左侧顶点数归纳证明：若有非空真子集 $S$ 满足等号，则分别匹配 $S$ 与 $N(S)$；在剩余图中，对任何员工集 $T$，原条件应用于 $S\cup T$ 给出 $|N(T)\setminus N(S)|\ge|T|$，所以剩余部分也可归纳匹配。
若每个非空真子集都满足严格不等式，任取一条边并删去其两个端点。剩余任意员工集的邻居至多减少一个，仍满足 Hall 条件，故可归纳匹配，再补回所选边。单个顶点是直接成立的基例。
因此原图存在一个完美匹配，作为第一轮。删掉其边后每个顶点的度数都变为 $d-1$；对剩余正则图重复，直到度数为零，得到 $d$ 个互不重叠的完美匹配。每轮覆盖所有人和任务，总共恰使用全部 $nd$ 条允许边一次。
\insight 局部度数条件通过边计数推出 Hall 条件；正则性又在删除一轮之后保留，从而把一次匹配提升为完整分解。

\problem{CM02}{非互素模数的同余构造}{广义中国剩余定理；扩展欧几里得}
给定正整数 $m,n$ 及整数 $a,b$，求同余组 $x\equiv a\pmod m$、$x\equiv b\pmod n$ 有解的充要条件，并给出可执行构造及解的周期。再求
\[
 x\equiv2\pmod6,\qquad x\equiv5\pmod9,
 \qquad x\equiv4\pmod5
\]
的全部整数解。
\solution 令 $g=\gcd(m,n)$。若有解，两式相减说明 $g\mid b-a$，这是必要条件。满足时写 $x=a+mt$，第二式等价于
\[
 \frac m g t\equiv\frac{b-a}{g}\pmod{n/g}.
\]
约分后的两个模系数互素，所以扩展欧几里得算法可求 $m/g$ 在模 $n/g$ 下的逆元，由此解出 $t$。若 $n/g=1$，则没有剩余限制，可取 $t=0$。因而条件也充分。
任意两个解的差同时被 $m,n$ 整除，故被 $\operatorname{lcm}(m,n)=mn/g$ 整除；反过来加上其任意整数倍仍是解。因此解恰为一个模 $mn/g$ 的剩余类。
例子中先写 $x=2+6t$，得到 $6t\equiv3\pmod9$，即 $2t\equiv1\pmod3$，所以 $t\equiv2\pmod3$、$x\equiv14\pmod{18}$。再写 $x=14+18k$，第三式给出 $3k\equiv0\pmod5$，即 $k\equiv0\pmod5$。全部解为
\[
 \boxed{x=14+90j,\qquad j\in\mathbb Z}.
\]
\insight 非互素模数首先需要检查相容性；约去公因子以后才能求逆，不能直接把不可逆的系数除掉。

\problem{CM03}{十二球的预定称量}{信息下界；三元编码}
十二枚外观相同的球中恰有一枚质量异常，可能较重或较轻，正常球质量相同。用无砝码天平，每次两盘放相同枚数的球，设计三次称量，且三次安排须事先固定、不随结果改变，以确定异常球及轻重。证明两次不够，并证明同类三次预定方案不能处理十三枚球。
\solution 每次结果为左重、平衡、右重，分别记 $1,0,-1$。十二球轻重共 $24$ 种状态，而两次只有 $3^2=9$ 种结果，所以至少三次。
以下矩阵每列对应一枚球，每行对应一次称量；$1$ 放左盘、$-1$ 放右盘、$0$ 不称：
\[
 \begin{array}{c|rrrrrrrrrrrr}
 &1&2&3&4&5&6&7&8&9&10&11&12\\\hline
 1&0&0&0&0&-1&-1&-1&1&1&1&-1&1\\
 2&0&-1&-1&-1&1&1&1&0&0&0&-1&1\\
 3&-1&1&0&-1&1&0&-1&-1&0&1&1&0
 \end{array}
\]
每行正负项各四个，所以称量合法。每列非零，且任意两列既不相同也不互为相反数。因此较重的第 $j$ 球产生第 $j$ 列向量，较轻时产生其相反向量；观察三次结果后即可唯一译码，所列 $24$ 个结果没有冲突。
若要处理十三球，必须用完 $3^3-1=26$ 个非零结果，每列与其相反数构成一对，十三列恰各代表一对。固定任意一行，在全部 $26$ 个非零三元向量中，该坐标非零的有 $2\cdot3^2=18$ 个，所以十三个代表列中这一行恰有九个非零项。九个球不可能均分到两盘，违背每次两盘球数相等。因此十三球不可能，十二球方案已达到这一类预定称量的上限。
\insight 信息计数只是必要条件；编码还必须满足实际称量的平衡约束。把轻重对应到互为相反的结果向量，可以把两者同时检查。
